% 10-validation.tex — Réalisation 4/4 : validation technique et physique (≈ 3 pages).

\chapter{Validation des simulations}
\label{chap:validation}

\section{Une chaîne de vérifications}

\begin{arediger}{1 page}
Synthèse des trois niveaux (détaillés dans la section~\ref{sec:tests}) :
tests unitaires de la physique contre des valeurs calculées à la main
(134 tests au moment de la rédaction, dont 53 pour le looping avec wagon et
son jeu) ; tests d'intégration Playwright sur neuf profils d'appareils
(tableau~\ref{tab:appareils}) ; vérification physique par un agent
indépendant qui redérive les formules et compare numériquement. Tableau :
niveau, outil, ce qui est vérifié, nombre de cas.
\end{arediger}

\section{Ce que la vérification indépendante a trouvé}

\begin{arediger}{1,5 page}
Exemple concret pour le jeu du seau : l'agent de vérification a confirmé la
physique du wagon (erreurs \(< 10^{-14}\) sur les formes fermées, \(< 2 \cdot 10^{-6}\)
sur les temps RK4) et relevé six écarts dans le jeu, corrigés ensuite : poids
de la barre sous-estimé d'un tiers, tension lue au mauvais bout de la corde,
dissipation numérique non signalée, durée trop courte pour de très petites
hauteurs, libellés. Montrer que ce dispositif attrape des erreurs qu'une
relecture humaine laisserait passer, et discuter ses limites.
\end{arediger}

\section{Performance et compatibilité}

\begin{arediger}{½ page}
Temps de calcul (trajectoire RK4 du wagon : quelques millisecondes ; essai du
seau : \SI{16}{ms} ; génération d'un défi : \SI{60}{ms}), fonctionnement en
\fichier{file://} sur Chromium, WebKit et Firefox, mise en page d'un seul écran
à partir de \(1000 \times 760\) pixels.
\end{arediger}
